% Authority: exact-scope leaf 14; full PDF page 93; printed p.76.
% U:p076.r28.par.005 | continuation; R28 par.005 visually replayed
\par\noindent
It follows from the theorem stated at the end of the preceding paragraph that, in order to decide
whether $-b$ is or is not a biquadratic residue with respect to $p$, everything reduces to the
question of knowing whether $t$ is or is not a quadratic residue with respect to $b$. We can now
show, by means of the results just obtained, that this latter question can be decided without
knowing $t$, that is, without its being necessary to put $p$ under the form $t^2-bu^2$. Suppose
first that $\varphi$ is not divisible by $b$. The relations $(\delta)$, which hold in this case,
can be presented in this way, observing that $b$ is a prime number $4n+3$:
% U:p076.r28.eq.006 | continuation; R28 eq.006 visually replayed
\[
\leg{\psi+t}{b}=1,\qquad \leg{\psi-t}{b}=-1.
\]
% U:p076.r28.par.006 | continuation; R28 par.006 visually replayed
\par\noindent
One immediately draws from the equation $p=t^2-bu^2$: $t^2\equiv p\md{b}$, a result which shows
that, on solving the congruence $\chi^2\equiv p\md{b}$, and denoting by $\chi$ one of its roots
taken at random, one will have $t\equiv\chi$, or $t\equiv-\chi\md{b}$. In the first case, one has:
% U:p076.r28.eq.007 | continuation; R28 eq.007 visually replayed
\[
\leg{t}{b}=\leg{\chi}{b},\qquad \leg{\chi+\psi}{b}=1,
\]
for it is evident that the expressions $\leg{t}{b}$ and $\leg{t+\psi}{b}$ do not change when one
puts, in place of $t$, another number $\chi$ which differs from $t$ only by a multiple of $b$.
Multiplying the last relations together, one obtains:
% U:p076.r28.eq.008 | continuation; R28 eq.008 visually replayed
\[
\leg{t}{b}=\leg{\chi(\chi+\psi)}{b}.
\]
% U:p076.r28.par.007 | continuation; R28 par.007 visually replayed
\par\noindent
Consider now the other case, in which $t\equiv-\chi\md{b}$. One will then have:
% U:p076.r28.eq.009 | continuation; R28 eq.009 visually replayed
\[
\leg{t}{b}=-\leg{\chi}{b},\qquad \leg{\psi+\chi}{b}=-1,
\]
the last of these relations resulting from the formula $\leg{\psi-t}{b}=-1$, if one puts, in place
of $-t$, the number $\chi$, which differs from it only by a multiple of $b$. Multiplying the
preceding relations together, one will find, as above:
% U:p076.r28.eq.010 | continuation; R28 eq.010 visually replayed
\[
\leg{t}{b}=\leg{\chi(\chi+\psi)}{b}.
\]
% U:p076.r28.par.008 | continuation; R28 par.008 visually replayed
\par\noindent
We have therefore arrived at this remarkable result: If $\varphi$ is not divisible by $b$, one can
decide very simply whether one has:
% U:p076.r28.eq.011 | continuation; R28 eq.011 visually replayed
\[
\leg{t}{b}=1\quad\hbox{or}\quad \leg{t}{b}=-1.
\]
